\section*{Aermacchi MB-326}

The Aermacchi MB-326 was designed as a trainer but then developed into a light attack aircraft. It featured tandem seating, a single Viper 11 turbojet engine, a low wing with wing-tip fuel tanks, and excellent performance for an aircraft of its class and era.

The initial MB-326 trainer (with no letter suffix) entered service with the Italian AMI in 1962. It had no provision for armament.

The MB-326B/E/F/H are the first combined trainer and attack versions. They are largely similar and all feature six under-wing weapon stations for up to 2000 lb of stores, including gun pods, bombs, and rocket pods. Versions after the B can also carry under-wing fuel tanks. The B was delivered to the Tunisian air force in 1965, the E to the Italian AMI for weapons training in 1968, the F to the Ghanian air force in 1965, and the H to the Australian RAAF and RAN starting in 1967. The H was also manufactured in Australia by the Commonwealth Aircraft Corporation (CAC) as the CA-30. 

The later MB-326G light attack version features a more powerful Viper 20-540 engine and a strengthened wing, allowing up to 4000 lb of stores to be carried. The GB version was delivered to the Argentine naval COAN starting in 1969, the Zairean air force in 1969, and Zambian air force in 1971. The GC version was delivered to the the Brazilian air force starting in 1971. The GC was also manufactured under license by Embraer as the EMB-326. In Brazilian service, the GC was known as the AT-26 or RT-26 Xavante. A number of Brazilian GCs were transferred to the Argentine navy after 1982 to replace losses in the South Atlantic War.

The MB-326M is largely similar to the G. It was manufactured both by Aermacchi and under license in South Africa by the Atlas Aircraft Corporation. Both versions served with the South African SAAF from 1966 as the Impala Mk I, and saw extensive combat during the South African Border War.

The MB-326K is a dedicated light attack aircraft and dispensed with the second crew member. It has two internal 30 mm DEFA cannon mounted in the fuselage, an armored, single-seat cockpit, an uprated Viper 632-43 engine, six stations for up to 4000 lb of stores, including gun pods, bombs, rocket pods, fuel tanks, four-camera photo-reconnaissance pod, and R.550 Magic infrared-homing missiles. It was used by the South African SAAF starting sometime after 1971. The K was also manufactured under license in South Africa from 1974 by the Atlas Aircraft Corporation, albeit with the earlier Viper 20-540 motor. Both variants of the K served in the SAAF as the Impala Mk.II.

The MB-326L was a two-seater version of the K. It was used by the air forces of Ghana, Tunisia, UAE, and Zaire from 1975. It was further developed into the MB-336.

ADCs are provided for the:
\begin{itemize}
\item MB-326B
\item MB-326E
\item MB-326F
\item MB-326G
\item MB-326H
\item MB-326K
\item MB-326L
\item MB-326M
\end{itemize}

See also:
\begin{itemize}
    \item Aermacchi MB-336
    \item Atlas Impala
    \item CAC CA-30
    \item Embraer AT-26.
\end{itemize}
